% from ega4-2-fr.tex
\oldpage[IV-2]{169}
de $X_\lambda$ au-dessus de $\mathfrak p$ appartiennent à
$\operatorname{Nor}(X_\lambda)$; en d'autres termes on a
$\bigcap_\lambda T_\lambda=\varnothing$; comme $Y$ est noethérien et les
$T_\lambda$ fermés, il existe un $\lambda$ tel que
$T_\lambda=\varnothing$; donc $B_\lambda$ est intégralement clos, et comme
son corps des fractions est $K'$, on a $B_\lambda=A'$. C.Q.F.D.
\end{proof}

\begin{proposition}[6.13.7]
\label{IV.2.s6.13.7-fr}
Soient $A$ un anneau noethérien, $X=\operatorname{Spec}(A)$. Les conditions
suivantes sont équivalentes:
\begin{enumerate}
\item[a)] Pour tout préschéma $X'$ localement de type fini sur $X$,
$\operatorname{Nor}(X')$ est ouvert dans $X'$;
\item[b)] Pour toute $A$-algèbre finie et intègre $A'$,
$\operatorname{Nor}(\operatorname{Spec}(A'))$ est ouvert dans
$\operatorname{Spec}(A')$;
\item[c)] Pour tout idéal premier $\mathfrak p$ de $A$ et toute extension
finie radicielle $K'$ du corps des fractions $K$ de $A/\mathfrak p$, il
existe une sous-$A$-algèbre $A'$ de $K'$, finie sur $A$, ayant $K'$ pour
corps des fractions et telle que $\operatorname{Nor}(\operatorname{Spec}(A'))$
soit ouvert dans $\operatorname{Spec}(A')$.
\end{enumerate}
\end{proposition}

\begin{proof}
Le fait que a) implique b) et que b) implique c) se prouve comme dans
(6.12.4). Pour montrer que c) entraîne a), compte tenu de (6.13.2), on se
ramène comme dans (6.12.4) à prouver que si $A'$ est une $A$-algèbre intègre
et de type fini, le point générique de $\operatorname{Spec}(A')$ est
intérieur à $\operatorname{Nor}(\operatorname{Spec}(A'))$. On distingue
alors deux cas comme dans (6.12.4) en prouvant d'abord le

\begin{lemma}[6.13.7.1]
\label{IV.2.s6.13.7.1-fr}
Soient $A$ un anneau intègre noethérien, $A'$ une $A$-algèbre intègre de
type fini, contenant $A$ et telle que le corps des fractions $K'$ de $A'$
soit une extension séparable du corps des fractions $K$ de $A$. Si
$\operatorname{Nor}(\operatorname{Spec}(A))$ est ouvert dans
$\operatorname{Spec}(A)$, $\operatorname{Nor}(\operatorname{Spec}(A'))$ est
ouvert dans $\operatorname{Spec}(A')$.
\end{lemma}

\begin{proof}
Il s'agit de prouver (compte tenu de (6.13.2)) que le point générique de
$\operatorname{Spec}(A')$ est intérieur à
$\operatorname{Nor}(\operatorname{Spec}(A'))$. La démonstration suit la
même marche que celle de (6.12.4.1), dont nous conservons les notations. On
remarque d'abord que l'on peut supposer que $A$ est intégralement clos, et
alors on sait que $A_1=A[t_1,\ldots,t_n]$ est intégralement clos (Bourbaki,
Alg. comm., chap. V, § 1, n$^o$ 3, cor. 2 de la prop. 13); on se ramène
ensuite au cas où $A'$ est un $A$-module libre de type fini; le raisonnement
de (6.12.4.1) prouve alors (en remplaçant au besoin $A$ par un anneau $A_f$
avec $f\neq0$) que les fibres $g^{-1}(z)$ du morphisme $g:X'\to X$ sont
régulières et a fortiori normales. En outre $g$ est plat et $X$ est normal,
donc (6.5.4, (ii)) $X'$ est normal.
\end{proof}

Ce lemme étant démontré, on passe au cas général comme dans (6.12.4, II)),
dont nous conservons encore les notations; appliquant l'hypothèse c), on
voit cette fois que $\operatorname{Nor}(X'_1)$ est ouvert et on se ramène
ainsi au cas où $X'_1$ est normal et $g:X'_1\to X'$ plat et surjectif; on
conclut cette fois que $X'$ est normal à l'aide de (6.5.4, (i)).
\end{proof}

\subsection{Changement de base et clôture intégrale}
\label{IV.2.s6.14-fr}

\begin{proposition}[6.14.1]
\EGAReaderTarget{ega4.c04.prop.6-14-1}{IV.2.s6.14.1-fr}
Soient $Y$, $Y'$ deux préschémas localement noethériens,
$g:Y'\to Y$ un morphisme normal (6.8.1). Alors, pour tout $Y$-préschéma
normal $X$, le préschéma $X'=X\times_Y Y'$ est normal.
\end{proposition}

Notons que dans cet énoncé on ne suppose pas $X$ localement noethérien.

\begin{lemma}[6.14.1.1]
\label{IV.2.s6.14.1.1-fr}
Soit $R$ un anneau composé direct d'un nombre fini de corps.
\begin{enumerate}
\item[(i)] Pour qu'un sous-anneau $A$ de $R$ ayant $R$ pour anneau total
des fractions soit normal, il faut et il suffit qu'il soit intégralement
fermé dans $R$.
\item[(ii)] Soit $(A_\lambda)$ une famille de sous-anneaux normaux de $R$;
si $A=\bigcap_\lambda A_\lambda$ admet $R$ pour anneau total des fractions,
$A$ est normal.
\end{enumerate}
\end{lemma}

\begin{proof}
(i) Comme $A$ est un sous-anneau de $R$, $A_\mathfrak p$ est un sous-anneau
de $R_\mathfrak p$ pour tout idéal premier $\mathfrak p$ de $A$, et
$R_\mathfrak p$ est un anneau de fractions de $A_\mathfrak p$; en outre
$R_\mathfrak p$ est nécessairement composé direct d'un nombre fini de
corps, donc tout élément non diviseur de $0$ dans $R_\mathfrak p$ est
inversible; cela prouve que $R_\mathfrak p$ est l'anneau total des fractions
de $A_\mathfrak p$. Si $A$ est intégralement fermé dans $R$,
$A_\mathfrak p$ est donc intégralement fermé dans $R_\mathfrak p$; mais si
$R_\mathfrak p$ est composé direct de 2 corps au moins, la fermeture
intégrale de $A_\mathfrak p$ dans $R_\mathfrak p$ est composée directe de
2 anneaux au moins non réduits à $0$, ce qui est absurde puisque
$A_\mathfrak p$ est un anneau local; donc $R_\mathfrak p$ est
nécessairement un corps et $A_\mathfrak p$ est intègre et intégralement clos,
ce qui par définition signifie que $A$ est normal. Inversement, si $A$ est
normal, $R_\mathfrak p$, anneau total des fractions d'un anneau intègre
$A_\mathfrak p$, est un corps et $A_\mathfrak p$ est intégralement fermé
dans $R_\mathfrak p$; si $x\in R$ est un élément entier sur $A$, son image
dans chaque $R_\mathfrak p$ est entière sur $A_\mathfrak p$, donc appartient
à $A_\mathfrak p$; on en conclut que $x\in A$ (Bourbaki, Alg. comm., chap.
II, § 3, n$^o$ 3, cor. 1 du th. 1), et par suite $A$ est intégralement fermé
dans $R$.

